\section{Probability Space}

\subsection{Probability Space}

\subsubsection{Random Experiment}

\begin{itemize}
\item A \textbf{random experiment} is an experiment in which the outcome varies in an unpredictable fashion when the experiment is repeated under the same conditions, i.e., \textit{random}.
\item An \textbf{outcome} of a random experiment is a result that cannot be decomposed into other results.
\end{itemize}

\subsubsection{Sample Space}

\begin{itemize}
\item The set of all possible outcomes for a random experiment is called the \textbf{sample space} $\Omega$.
\end{itemize}

\subparagraph{Examples:}

\begin{enumerate}[resume]
\item Tossing a coin: $\Omega=\{H, T\}$
\item Tossing a coin twice: $\Omega=\{H H, H T, T H, T T\}$
\item Taking EEE 5544 and getting a grade: $\Omega=\{A, A -, B+, B, B -, C+, C, D, F\}$
\item Picking an integer at random: $\Omega=\{\ldots, -3, -2, -1,0,1,2,3, \ldots \}$
\item Checking the temperature (in $K$) at a random location $\Omega=(0, \infty)$
\end{enumerate}

\begin{itemize}
\item \textbf{Finite} $\Omega$: Examples 1, 2, 3
\item \textbf{Countably infinite} $\Omega$: Example 4
\item \textbf{Discrete} $\Omega$: Examples 1--4
\item \textbf{Continuous (uncountably infinite)} $\Omega$: Example 5
\end{itemize}

\subsubsection{Event Class}

\begin{itemize}
\item An \textbf{event} is a set of outcomes (a subset of $\Omega$) that we can ``quantify its likelihood'' (see below).
\item Two special events:

\begin{itemize}
\item $\Omega =$ \textbf{certain event}
\item $\emptyset =$ \textbf{impossible event}
\end{itemize}


\item An \textbf{event class} $\mathcal{F}$ of a random experiment is a collection of events.
\end{itemize}

\paragraph{Discrete sample space}

\begin{itemize}
\item For a discrete $\Omega$, we typically consider the event class that is the collection of all subsets of $\Omega$, i.e., $\mathcal{F} = 2^{\Omega}$, which is called the \textit{power set} of $\Omega$.
\end{itemize}

\subparagraph{Examples:}

\begin{enumerate}[resume]
\item Tossing a coin: $\mathcal{F} =\{\emptyset,\{H\},\{T\},\{H T\}\}$.
\item Tossing a coin twice:

\begin{equation*}
\mathcal{F} = \left\{\emptyset,\{H H\},\{T T\},\{H T\},\{T H\},
 \{H H, T T\},\{H H, H T\}, \ldots, 
\{H H, T T, H T\}, \ldots,\{H H, T T, H T, T H\} \right\}.
\end{equation*}


\item If $\Omega=\left\{a_{1}, a_{2}, \ldots, a_{N}\right\}$, then

\begin{equation*}
\mathcal{F}=\left\{\emptyset,\{a_{1}\},
\ldots,\{a_{N}\}, \{a_{1}, a_{2}\},
\ldots,\{a_{1}, a_{2}, a_{3}\}, \ldots,
\{a_{1}, a_{2}, \ldots, a_{N}\}\right\}.
\end{equation*}

Altogether $2^{N}$ events (why?).


\item For a countably infinite $\Omega = \left\{a_{1}, a_{2},
\ldots \right\}$, we can use the same idea as in the finite case using the power set as the event space, i.e.,

\begin{equation*}
\mathcal{F} = 2^{\Omega} = \left\{\emptyset,\{a_{1}\},
\ldots,\{a_{N}\}, \{a_{1}, a_{2}\},
\ldots,\{a_{1}, a_{2}, a_{3}\}, \ldots,
\{a_{1}, a_{2}, \ldots, a_{n}\}, \ldots \right\}
\end{equation*}

by enumerating all subsets of $\Omega$. However, we need a better way to describe the power set.
\end{enumerate}

\paragraph{Continuous sample space}

\begin{itemize}
\item For a continuous sample space, e.g., $\Omega = (0,\infty)$, there are uncountably infinite number of outcomes and hence the method of enumeration used in the finite case can't be used to construct the event class anymore.
\item Luckily mathematicians have solved the problem for us. To discuss the solution, we need to study a bit about set operations.
\end{itemize}

\subsubsection{Set operations}

\begin{itemize}
\item \textbf{Union:} $A \cup B=\{x \in \Omega: x \in A$ or $x \in B\}$
\item \textbf{Intersection:} $A \cap B=\{x \in \Omega: x \in A$ and $x \in B\}$
\item \textbf{Complement:} $A^{c}=\{x \in \Omega: x \notin A\}$
\item \textbf{Difference:} $A \setminus B=\{x \in \Omega: x \in A$ and $x \notin B\}=A \cap B^{c}$
\end{itemize}

\subparagraph*{\textit{Proof:}}

We give the proof of 1. as an example here:\newline
 (i) Since $x \in(A \cup B)^{c} \Rightarrow x \notin A \cup B \Rightarrow 
x \notin A$ and $x \notin B \Rightarrow x \in A^{c}$ and $x \in
B^{c} \Rightarrow x \in A^{c} \cap B^c$, $(A \cup B)^{c} \subseteq A^{c} \cap B^{c}$.

(ii) Since $x \in A^{c} \cap B^{c} \Rightarrow x \in A^{c}$ and $x \in B^{c} \Rightarrow x \notin A$ and $x \notin B
\Rightarrow x \notin A \cup B \Rightarrow x \in(A \cup B)^{c}$, $A^{c} \cap B^{c} \subseteq (A \cup B)^{c}$.

Combining (i) and (ii), we have $(A \cup B)^{c} = A^{c} \cap B^{c}$.

\subsubsection{Field (Algebra)}

\begin{itemize}
\item An \textbf{field (algebra)} is a collection of subsets of $\Omega$ that is closed under union and complement, and contains both $\emptyset$ and $\Omega$. That is, $\mathcal{F}$ is a field (algebra) if it satisfies the following three conditions:

\begin{itemize}
\item $\emptyset \in \mathcal{F}$ and $\Omega \in \mathcal{F}$.
\item If $A \in \mathcal{F}$ and $B \in \mathcal{F}$, then $A \cup B \in \mathcal{F}$.
\item If $A \in \mathcal{F}$, then $A^{c} \in \mathcal{F}$.
\end{itemize}


\item Operationally speaking, a field contains all subsets of $\Omega$ that can be obtained from applying a finite number of set operations.
\end{itemize}

\subsubsection{$\sigma$-Field ( $\sigma$-Algebra)}

\begin{itemize}
\item A field $\mathcal{F}$ is a $\sigma$-field if it satisfies the additional condition:

\begin{itemize}
\item If $A_{1}, A_{2}, \ldots \in \mathcal{F}$, then $\bigcup_{i=1}^{\infty} A_{i} \in \mathcal{F}$.
\end{itemize}


\item Operationally speaking, a $\sigma$-algebra contains all subsets of $\Omega$ that can be obtained from a countably number of set operations.
\end{itemize}

\paragraph{Generating $\sigma$-field}

\begin{itemize}
\item For a continuous (uncountably infinite) sample space $\Omega$, we often \textit{generate} $\sigma$-fields from some collections of subsets of $\Omega$.
\item For a collection $\mathcal{A}$ of subsets of $\Omega$, consider the $\sigma$-field generated by $\mathcal{A}$, i.e., the smallest $\sigma$-field containing all subsets of $\Omega$ that can be formed by performing countable set operations on subsets in $\mathcal{A}$.
\item The $\sigma$-field generated by $\mathcal{A}$ is denoted by $\sigma(\mathcal{A})$.
\end{itemize}

\subsubsection{Event Class for Continuous Sample Space}

\begin{itemize}
\item For a continuous (uncountably infinite) sample space $\Omega$, we will consider event classes that are $sigma$-fields of $\Omega$.
\end{itemize}

\subparagraph{Example: $\Omega=\mathbb{R}$}

\begin{itemize}
\item Consider the collection $\mathcal{A}$ of all the subsets of $\mathbb{R}$ in the form of $( -\infty, a)$ for some $a \in \mathbb{R}$.
\item The $\sigma$-field generated by $\mathcal{A}$, called the \textit{Borel $\sigma$-field}, is commonly used as an event class for the sample space $\mathbb{R}$.
\item The Borel $\sigma$-field contains essentially all events of practical interest.
\item While the \textit{Borel sets} (i.e., subsets in the Borel $\sigma$-field) form a very rich class of subsets of $\mathbb{R}$, there are subsets of $\mathbb{R}$ that are not Borel sets. This turns out to be critically important in allowing the ``likelihood'' of the events to be ``quantifiable'' as we discussed before.
\end{itemize}

\subsubsection{Description of Random Experiment}

\begin{itemize}
\item We can now use the sample space $\Omega$ and an event class $\mathcal{F}$ to mathematically describe a random experiment.
\item In addition, we need a way to ``quantify the likelihood of the events'' as allured to above. This is done by a mapping called \textit{probability measure $P$}, which will be discussed in the next section.
\item In summary, the math object that describes a random experiment is the triple $(\Omega, F, P)$ called the \textbf{probability space} of the random experiment.
\end{itemize}